\chapter{Multivariate, Isogeny-Based, and Other Approaches}
\label{ch:other-approaches}

\abstract*{Not every post-quantum proposal has survived scrutiny. This chapter examines the multivariate and isogeny-based families---including the spectacular collapse of SIKE in 2022---not merely as historical curiosities, but as case studies in what makes a cryptographic assumption trustworthy. Understanding how schemes break is as important as understanding how they work.}

% =============================================================
% SOURCE: notes.tex §3.4.4 (lines 4237–4488)
% REUSE: ~40%. Good narrative on failures. Expand isogeny math,
%        add CSIDH, draw lessons.
% =============================================================

\section{Multivariate Polynomial Cryptography}
\label{sec:multivariate}

% TODO: Adapt from notes.tex §3.4.4
% - The MQ problem: definition, NP-completeness
% - Oil-Vinegar construction
% - The pattern of breaks: C*, HFE, SFLASH, Rainbow
% - Rainbow's fall (Beullens 2022): the attack explained
% - Why structure keeps leaking: a fundamental design tension
% - UOV: the survivor — current status
% - Cite: Kipnis-Shamir 1999, Beullens 2022

\section{Isogeny-Based Cryptography}
\label{sec:isogenies}

% TODO: Expand significantly
% - Elliptic curves and isogenies: mathematical background
% - Supersingular isogeny graphs: the expander graph structure
% - SIDH/SIKE: the construction and its elegant key exchange
% - The Castryck-Decru attack (2022): exploiting auxiliary points
% - Why the break was possible: the torsion point information leak
% - Timeline: from NIST finalist to broken in a weekend
% - Cite: Jao-De Feo 2011, Castryck-Decru 2022

\section{Group Actions and CSIDH}
\label{sec:csidh}

% TODO: New section
% - Class group actions on elliptic curves
% - CSIDH: commutative supersingular isogeny DH
% - Why CSIDH survives the SIKE attack
% - Challenges: class group computation, constant-time implementation
% - Current status and open problems
% - Cite: Castryck et al. (CSIDH) 2018

\section{Lessons Learned}
\label{sec:lessons}

% TODO: New section (pedagogically valuable)
% - What makes a cryptographic assumption trustworthy?
% - The role of time: 45 years (McEliece) vs 10 years (SIKE)
% - Structured vs unstructured problems
% - The danger of auxiliary information
% - Diversification: why multiple families matter
% - Crypto-agility as a survival strategy

%% Exercises
\begin{prob}
\label{prob:ch11-mq}
Consider a system of 2 quadratic equations in 3 variables over $\F_2$. Enumerate all solutions by brute force. Then explain why this approach becomes infeasible as the number of variables grows, despite the problem being NP-complete.
\end{prob}

\begin{prob}
\label{prob:ch11-sike}
Explain the key exchange protocol of SIDH at a high level. Identify the auxiliary torsion point information that Alice sends to Bob. Why is this information necessary for the protocol, and how did Castryck and Decru exploit it?
\end{prob}

\begin{prob}
\label{prob:ch11-trust}
Propose a set of criteria for evaluating the trustworthiness of a post-quantum cryptographic assumption. Apply your criteria to: (a)~Module-LWE, (b)~the syndrome decoding problem, (c)~the MQ problem, (d)~the supersingular isogeny problem (post-SIKE). Rank them and justify your ranking.
\end{prob}
